\SourceLeaf{538}\par\medskip\Needspace{5\baselineskip}\textit{Nachlass.}
\begingroup\fontsize{10}{13}\selectfont
\begin{longtable}{@{}>{\raggedright\arraybackslash}p{.86\linewidth}@{\quad}r@{}}
\multicolumn{2}{@{}l}{Untersuchung der Reste:}\\[3pt]
Sechstes Kapitel. Erster Teil. Lösung der Kongruenz $x^m-1\equiv0$ .\,.\,.\,.\,.\,.\,.\,.\,. & Seite 199 \\[3pt]
Achtes Kapitel. Allgemeine Untersuchungen über Kongruenzen .\,.\,.\,.\,.\,.\,.\,.\,.\,. & — 212 \\[3pt]
Weiterentwicklung der Untersuchungen über reine Gleichungen .\,.\,.\,.\,.\,.\,.\,.\,.\,.\,.\,. & — 243 \\[3pt]
Beweis einiger Sätze über die Perioden der Klassen binärer quadratischer Formen .\,.\,.\,.\,.\,.\,.\,.\,.\,.\,.\,.\,.\,.\,.\,.\,.\,.\,.\,.\,.\,.\,.\,. & — 266 \\[3pt]
Über den Zusammenhang zwischen der Anzahl der Klassen, in die binäre quadratische Formen zerfallen, und ihrer Determinante. I. II\,.\,. X .\,.\,.\,.\,.\,.\,.\,.\,.\,.\,.\,.\,.\,.\,.\,.\,. & — 269 \\[3pt]
Geometrischer Aspekt der ternären Formen .\,.\,.\,.\,.\,.\,.\,.\,.\,.\,.\,.\,.\,.\,.\,.\,. & — 305 \\[3pt]
Zur Theorie der biquadratischen Reste. I.\,.\,. VI .\,.\,.\,.\,.\,.\,.\,.\,.\,.\,.\,.\,.\,. & — 313 \\[3pt]
Zur Theorie der komplexen Zahlen. I.\,.\,. VI .\,.\,.\,.\,.\,.\,.\,.\,.\,.\,.\,.\,.\,.\,. & — 387 \\[3pt]
Tafel des quadratischen Charakters der Primzahlen .\,.\,.\,.\,.\,.\,.\,.\,.\,.\,.\,.\,.\,. & — 399 \\[3pt]
Tafel zur Umwandlung gemeiner Brüche in Dezimalbrüche .\,.\,.\,.\,.\,.\,.\,.\,.\,.\,. & — 411 \\[3pt]
Tafel der Häufigkeit der Primzahlen .\,.\,.\,.\,.\,.\,.\,.\,.\,.\,.\,.\,.\,.\,.\,.\,.\,.\,.\,. & — 435 \\[3pt]
Tafel der Anzahl der Klassen binärer quadratischer Formen .\,.\,.\,.\,.\,.\,.\,.\,.\,.\,. & — 449 \\[3pt]
Tafel zur Kreisteilung .\,.\,.\,.\,.\,.\,.\,.\,.\,.\,.\,.\,.\,.\,.\,.\,.\,.\,.\,.\,.\,.\,.\,. & — 477 \\[3pt]
Neue Quadratur des Kreises .\,.\,.\,.\,.\,.\,.\,.\,.\,.\,.\,.\,.\,.\,.\,.\,-\,.\,.\,.\,.\,.\,.\,. & — . 497 \\[3pt]
Zur Berechnung der Logarithmen .\,.\,.\,.\,.\,.\,.\,.\,.\,.\,.\,.\,.\,.\,.\,.\,.\,.\,.\,. & — 501 \\[3pt]
Die ersten zehntausend Quadrate .\,.\,.\,.\,.\,.\,.\,.\,.\,.\,.\,.\,.\,.\,.\,.\,.\,.\,.\,.\,.\,.\,. & — 504 \\[3pt]
Indizes der Primzahlen im höheren Zahlenbereich .\,.\,.\,.\,.\,.\,.\,.\,.\,,\,.\,.\,.\,.\,. & — 506 \\[3pt]
Hilfstafel zur Lösung der Gleichung $A=fxx+gyy$ .\,.\,.\,.\,.\,.\,.\,.\,.\,.\,.\,. & — 509 \\[3pt]
Achter Abschnitt. Ausführlichere Betrachtung einiger Untersuchungen zur Kreisteilung.  & — 510 \\[3pt]
Briefe von \textsc{Gauss} an \textsc{Dirichlet} .\,.\,.\,.\,.\,.\,.\,.\,.\,.\,.\,.\,.\,.\,.\,.\kern.6em.\,.\,.\,. & — 514 \\[3pt]
\end{longtable}\endgroup
\begingroup\fontsize{10}{13}\selectfont
\begin{longtable}{@{}>{\raggedright\arraybackslash}p{.86\linewidth}@{\quad}r@{}}
\textit{Bemerkungen} .\,.\,.\,.\,.\,.\,.\,.\,.\,.\,.\,.\,.\,.\,.\,.\,.\,.\,.\,.\,.\,.\,.\,.\,.\,.\,.\,. & — 519 \\[3pt]
\end{longtable}\endgroup
\BookRule{.25\textwidth}
\begin{center}
GÖTTINGEN,\par\smallskip
DRUCK DER DIETRICHSCHEN UNIVERSITÄTSBUCHDRUCKEREI.\par\smallskip
W. FR. KAESTNER.
\end{center}
